% =====================================================================
%  Gaya Buku Pedoman modulstis
%  Warna mengikuti bps_palette() paket modulstis. Mesin: XeLaTeX.
% =====================================================================
\usepackage{xcolor}
\definecolor{bpsnavy}{HTML}{123A6B}
\definecolor{bpsbiru}{HTML}{0B6FB4}
\definecolor{bpsbirumuda}{HTML}{6FB3E0}
\definecolor{bpsoranye}{HTML}{F58220}
\definecolor{bpshijau}{HTML}{00A651}
\definecolor{abuteks}{HTML}{45505C}
\definecolor{merahsalah}{HTML}{C0392B}

\usepackage{tikz}
\usetikzlibrary{calc,positioning,arrows.meta,shapes.misc}
\usepackage{graphicx}
\usepackage{booktabs}
\usepackage{array}
\usepackage{enumitem}
\setlist{itemsep=2pt, topsep=4pt}
\usepackage{ragged2e}

% ---------- Pemenggalan kata bahasa Indonesia ----------
% Babel memuat pola 'indonesian' (hyph-id). Daftar ini menambah pemenggalan
% untuk istilah yang sering muncul di buku dan mencegah nama dipenggal.
\hyphenation{mo-dul mo-dul-stis spe-si-fi-ka-si pra-tik-um ke-lu-ar-an ren-der
  pe-nyu-sun ber-ka-s Win-dows Quar-to}

% ---------- Paragraf ----------
\setlength{\parindent}{0pt}
\setlength{\parskip}{5pt plus 1pt}
\raggedbottom
\setlength{\emergencystretch}{3em}

% ---------- Judul (KOMA-Script) ----------
\setkomafont{disposition}{\sffamily\bfseries\color{bpsnavy}}
\setkomafont{chapter}{\Huge}
\setkomafont{section}{\Large\color{bpsbiru}}
\setkomafont{subsection}{\large\color{bpsnavy}}
\setkomafont{subsubsection}{\normalsize\color{bpsnavy}}
\setkomafont{part}{\Huge\color{white}}
\setkomafont{partnumber}{\Large\color{white}}
\setkomafont{chapterprefix}{\large\color{bpsoranye}}
\RedeclareSectionCommand[beforeskip=0pt, afterskip=18pt]{chapter}
\RedeclareSectionCommand[beforeskip=14pt plus 2pt, afterskip=6pt]{section}
\RedeclareSectionCommand[beforeskip=10pt plus 2pt, afterskip=4pt]{subsection}
\renewcommand*{\chapterformat}{%
  \mbox{\colorbox{bpsnavy}{\textcolor{white}{\,\thechapter\,}}\enskip}}
\renewcommand*{\chapterlinesformat}[3]{%
  \parbox{\linewidth}{\raggedright #2#3\par
  \vspace{4pt}{\color{bpsoranye}\rule{\linewidth}{1.6pt}}}}

% Halaman bagian (part): latar biru tua penuh
\renewcommand*{\partpagestyle}{empty}
\renewcommand*{\partformat}{\partname~\thepart}
\addtokomafont{partentry}{\color{bpsnavy}}
\makeatletter
\renewcommand*{\partheadstartvskip}{%
  \begin{tikzpicture}[remember picture, overlay]
    \fill[bpsnavy] (current page.south west) rectangle (current page.north east);
    \fill[bpsbiru] ($(current page.south west)+(0,5cm)$) --
          ($(current page.south east)+(0,9cm)$) -- (current page.south east) --
          (current page.south west) -- cycle;
    \fill[bpsoranye] ($(current page.south west)+(0,4.4cm)$) rectangle
          ($(current page.south east)+(0,4.6cm)$);
  \end{tikzpicture}\vspace*{6cm}}
\makeatother

% ---------- Kepala dan kaki halaman ----------
\usepackage[headsepline=1pt]{scrlayer-scrpage}
\clearpairofpagestyles
\addtokomafont{pageheadfoot}{\small\color{abuteks}}
\addtokomafont{headsepline}{\color{bpsbiru}}
\ihead{\textcolor{bpsnavy}{\textbf{Pedoman \texttt{modulstis}}}}
\ohead{\leftmark}
\cfoot*{\colorbox{bpsnavy}{\textcolor{white}{\,\pagemark\,}}}
\automark[chapter]{chapter}
\renewcommand*{\chaptermarkformat}{Bab~\thechapter\enskip}

% ---------- Kotak (callout Quarto) ----------
% Warna kotak disetel ulang setelah Quarto mendefinisikannya.
\AtBeginDocument{%
  \definecolor{quarto-callout-tip-color}{HTML}{00A651}%
  \definecolor{quarto-callout-note-color}{HTML}{0B6FB4}%
  \definecolor{quarto-callout-warning-color}{HTML}{F58220}%
  \definecolor{quarto-callout-important-color}{HTML}{C0392B}%
  \definecolor{quarto-callout-caution-color}{HTML}{123A6B}%
  \definecolor{quarto-callout-tip-color-frame}{HTML}{00A651}%
  \definecolor{quarto-callout-note-color-frame}{HTML}{0B6FB4}%
  \definecolor{quarto-callout-warning-color-frame}{HTML}{F58220}%
  \definecolor{quarto-callout-important-color-frame}{HTML}{C0392B}%
  \definecolor{quarto-callout-caution-color-frame}{HTML}{123A6B}%
}

% ---------- Langkah bernomor ----------
% \langkah{judul} membuat penanda langkah berwarna oranye dengan nomor otomatis.
\newcounter{langkah}[chapter]
\newcommand{\langkah}[1]{%
  \par\addvspace{6pt}\refstepcounter{langkah}%
  \noindent\colorbox{bpsoranye}{\textcolor{white}{\sffamily\bfseries\,Langkah \thelangkah\,}}%
  \enskip{\sffamily\bfseries\color{bpsnavy}#1}\par\nopagebreak\vspace{2pt}}

% ---------- Penanda sistem operasi ----------
\newcommand{\osbadge}[2]{\colorbox{#1}{\textcolor{white}{\sffamily\bfseries\small\,#2\,}}}
\newcommand{\windows}{\osbadge{bpsbiru}{Windows}}
\newcommand{\macos}{\osbadge{abuteks}{macOS}}
\newcommand{\linux}{\osbadge{bpshijau}{Linux}}

% ---------- Gambar berbingkai ----------
\newcommand{\bingkai}[2][0.47\linewidth]{%
  \fcolorbox{bpsnavy!40}{white}{\includegraphics[width=#1]{#2}}}

% ---------- Baris kode panjang dipenggal, tidak melewati tepi kotak ----------
\usepackage{fvextra}
\DefineVerbatimEnvironment{Highlighting}{Verbatim}{breaklines,breakanywhere,commandchars=\\\{\}}
\fvset{breaklines=true, breakanywhere=true}
